\documentclass[12pt]{article}
\usepackage[utf8]{inputenc}
\usepackage[T1]{fontenc}
\usepackage[margin=1in]{geometry}
\usepackage{setspace}
\usepackage{xcolor}
\usepackage[colorlinks=true, urlcolor=blue!60!black, linkcolor=blue!60!black]{hyperref}
\setlength{\parindent}{0em}
\setlength{\parskip}{0.6em}
\renewcommand{\baselinestretch}{1.0}
\begin{document}
\begin{center}
{\LARGE \textit{SEC v. Jarkesy}, 603 U.S. 109 (2024)\par}
\vspace{0.5em}
{\large Case Brief --- Two-Page Edition\par}
\end{center}

\section{Memory Jogger}
SEC civil penalties for securities fraud require a jury trial in an Article~III court; the public rights exception does not apply (\href{https://supreme.justia.com/cases/federal/us/603/22-859/}{603 U.S.\ 109}).

\section{Detailed Case Facts}
Dodd-Frank let the SEC seek civil penalties either in federal court, with a jury, or in its own administrative proceedings, without one. The SEC brought an in-house fraud action against George Jarkesy and his adviser Patriot28 for misleading investors in two hedge funds, and imposed a \$300,000 penalty and other sanctions.

\section{Procedural History}
A divided Fifth Circuit vacated the order on Seventh Amendment, nondelegation, and removal grounds. The Supreme Court granted certiorari.

\section{Judicial Votes}
Affirmed, 6--3. Roberts wrote for Thomas, Alito, Gorsuch, Kavanaugh, and Barrett; Gorsuch, joined by Thomas, concurred; Sotomayor, joined by Kagan and Jackson, dissented.

\section{Holding}
A statutory claim that is ``legal in nature'' carries the jury right, and punitive civil penalties make that showing ``all but dispositive.'' Securities fraud resembles common-law fraud and is not a public right, so Congress may not assign it to agency adjudication. Other grounds were not reached.

\section{Analysis of Opinions}
Roberts held that the penalties punish and deter rather than compensate, that securities fraud borrows common-law terms bringing their ``old soil,'' and that the public rights exception covers only historic categories like revenue, customs, and immigration. \textit{Atlas Roofing} was distinguished as involving duties ``unknown to the common law.'' Gorsuch tied the jury right to Article~III and due process. Sotomayor called the ruling a ``power grab'' that upends two centuries of precedent and threatens over 200 statutes. \textit{FCC v.\ AT\&T} (2026) later upheld FCC forfeitures enforceable only after a jury trial.

\section{Examples: Future Applications}
\textbf{Same-Side:} in-house penalties for commodities fraud require a jury. \textbf{Opposite-Side:} administrative fines on carriers under the immigration laws do not (\textit{Stranahan}). \textbf{Fence-Sitter:} immediately enforceable EPA penalties for technical emission limits.

\section{Critique}
Critics say the ``all but dispositive'' remedy test leaves \textit{Atlas Roofing} on razor-thin footing and offers no clear public-rights test. Proponents praise a clear baseline barring binding executive fines without a jury.

\section{Key Quotations}
\begin{quote}
``In this case, the remedy is all but dispositive. \ldots\ In sum, the civil penalties in this case are designed to punish and deter, not to compensate. They are therefore `a type of remedy at common law that could only be enforced in courts of law.' ''
\end{quote}
\noindent --- Chief Justice Roberts, Majority Opinion, 603 U.S.\ at 123, 125.

\end{document}
